% Section IV: Demonstration collection and policy deployment (target 1.10 pages with Table II).
% Sources: docs/data_rules.md, docs/experiments.md, docs/decisions.md (rec-1, rec-6, p3-a, p3-b, t3-a .. t3-e),
% session folders data/demos/real_* (episodes.json, notes.txt), configs/policy.yaml, checkpoint config.json.
% Evaluation placements, trials and success criteria: Section V.
\section{Demonstration Collection and Policy Deployment}
\label{sec:workflow}

\subsection{Task and Demonstration Protocol}
\label{sec:protocol}
One episode packs all three items in a fixed order: the carrot into the black
box, then the Oreo into the large hole and the RITZ into the small hole of the
meal box. A fixed order gives one correct action at every moment, whereas a
free order would let the policy average between choices. Before each episode,
the items are placed at random positions within the view of the top camera,
not touching each other or the containers. Which item lies nearest the robot
also varies, so the order must be inferred from the item types rather than from
their positions. Each item is oriented so that the fingers close across it
without wrist rotation. The operator starts each episode with the gripper
open, from a different pose within about 10~cm of the ready pose.

Every item is handled with the same sequence: approach above it, descend
vertically, close, lift, move above the target, lower slightly, open, and
lift. After the last item, the operator brings the gripper back above the
item area, near the ready pose, and ends the episode there; this pose is also
where the next episode starts. The operator moves steadily, without hesitation or searching. An episode is
discarded when an item drops or lands outside its target (including the wrong
hole), when anything is hit, when the motion pauses for more than about 2~s,
or when a camera reports an error. A natural re-grasp after a slip is kept
if the episode still ends correctly. Episodes end with a keyboard key or at a time
limit (180~s for three items), and every session is checked with the converter before the next one
begins.

\emph{Recovery demonstrations.}
A pilot policy that saw only successful grasps closed beside an item, lifted
the empty gripper, and carried it to the box. The gripper position cannot
reveal the miss, because the soft items and an empty grasp both leave the
fingers nearly closed, and the motor current rises in only some of the real
grasps (Section~\ref{sec:interface}); only the grip camera shows the miss
reliably. Recovery episodes
therefore contain one deliberate miss: the operator closes beside an item,
lifts the empty gripper, pauses for 0.5--1~s so that the grip camera sees the
empty fingers, opens, moves directly above the item, and continues as in a
normal episode. In the three-item recovery set, each item is missed eight times,
covering four directions on the table. The miss lies 2.5--15.6~cm from the
following grasp (median 8.5~cm, center to center), and the empty lift is
8--15~cm.

\emph{Collection stages.}
Data were collected in stages, so that each policy's failures could motivate
the next set (Table~\ref{tab:datasets}). The first stages placed the carrot
and one snack per episode, normal (2N) and with recovery (2R). A policy trained
on them had no demonstrated action once both snacks lay on the table, so all
later episodes contain all three items: normal (3N) and recovery (3R). The
earlier episodes remain valid under the same rule as scenes with one snack
absent. The evaluated policy is trained on the union of the four sets,
\emph{all} (140 episodes). Of 66 attempts in the three-item normal session, 51
were kept, and 7 of these include a natural re-grasp. Cameras, lighting,
container positions, and teleoperation settings were the same in every session.

\begin{table}[t]
\caption{Demonstration Sets and Training Datasets}
\label{tab:datasets}
\centering
\footnotesize
\setlength{\tabcolsep}{4pt}
\begin{tabular}{@{}llrrr@{}}
\toprule
Set & Content & Episodes & Frames & Mean (s)\\
\midrule
2N & Carrot + one snack, normal & 49 & 70{,}550 & 48.0\\
2R & Carrot + one snack, recovery & 16 & 24{,}862 & 51.8\\
3N & Three items, normal & 51 & 89{,}899 & 58.8\\
3R & Three items, recovery & 24 & 43{,}272 & 60.1\\
\midrule
all & 2N + 2R + 3N + 3R & 140 & 228{,}583 & --\\
\bottomrule
\end{tabular}
% Frames at 30 Hz; mean = frames / episodes / 30. Names in docs/experiments.md: 2N = u1, 2R = u2,
% 3N = u4n, 3R = u4r, all = u4 (u4a = u1 + u2 + u4n, an intermediate training set, not evaluated).
\end{table}

\subsection{Observation, Action, and Dataset Preparation}
\label{sec:dataset}
The state and the action at time step $t$ are
\begin{align}
 \mathbf{s}_t &= [q_{1,t},\ldots,q_{6,t},\,g^{\mathrm{pos}}_t],\\
 \mathbf{a}_t &= [x^{\mathrm{cmd}}_t,\,y^{\mathrm{cmd}}_t,\,z^{\mathrm{cmd}}_t,\,g^{\mathrm{cmd}}_t],
\end{align}
where $q_{1..6}$ are the measured follower joint angles in degrees, $g^{\mathrm{pos}}$
is the measured gripper opening, $(x,y,z)^{\mathrm{cmd}}$ is the shaped TCP
target in millimeters in the base frame (Section~\ref{sec:limits}), and
$g^{\mathrm{cmd}}$ is the trigger command. Both gripper values range from zero
(closed) to one (open). The label is thus the target that the robot was
commanded to reach, not its measured motion; the orientation follows from
the position by the rule of Section~\ref{sec:interface} and is not learned.
The observation adds the three camera images. Because the joint angles are
part of the observation, the policy is specific to the RB5.

A converter builds the dataset from the raw logs on a 30~Hz grid starting at
each episode's start. Targets and joint states, recorded at 50~Hz, are
interpolated to the grid times. For each camera, the frame captured nearest to
each grid time is used; in the four collection sessions, this frame lay within
17.6~ms of the grid time for 99\% of the samples. Episodes in which the clutch
was engaged for less than 1~s or the target moved less than 30~mm are left out,
as are episodes with camera gaps. The images are reduced from $640\times480$
to $320\times240$ with the same resampling function used at run time. The
result is stored in the LeRobot dataset format~\cite{cadene2024lerobot}, with
images as compressed video. Datasets from separate sessions are merged without
conversion, and each dataset keeps a copy of the settings used to build it.
% TODO: add the LeRobot entry (cadene2024lerobot) to references.bib.

\subsection{ACT Training and Autonomous Execution}
\label{sec:execution}
The policy is ACT as implemented in LeRobot~0.6.1. A ResNet-18 backbone
initialized from ImageNet encodes each camera image, and a transformer with
a 512-dimensional model and a conditional variational autoencoder (KL weight
10) predicts a chunk of 100 actions (3.3~s) from one observation. States,
actions, and images are normalized by their dataset mean and standard
deviation. Training uses a batch size of 8, a constant learning rate of
$10^{-5}$, and 100{,}000 steps, about 3.5 passes over \emph{all}. On the host's
laptop GPU this takes about 7~h. The checkpoint at the final step is
evaluated; no checkpoint is selected by test results.

At run time, the policy process predicts a new chunk from the latest images
and robot state whenever the previous inference has finished, at most once
per 30~Hz step. Each chunk carries the capture time of its observation. In
every 20~ms servo cycle, each recent chunk that covers the current time
contributes its action, interpolated between its 30~Hz steps, and the
contributions are averaged as in ACT~\cite{zhao2023act}:
\begin{equation}
\mathbf{a}(t) = \frac{\sum_{i} w_i\,\mathbf{a}^{(i)}(t)}{\sum_{i} w_i},\qquad
w_i = e^{-m i},
\label{eq:ensemble}
\end{equation}
where $i=0$ denotes the oldest chunk and $m=0.01$. Because the averaging is
done by time, a late or skipped inference only removes one chunk from the
average. The position passes through the command shaping of
Section~\ref{sec:limits}, and the gripper value goes to the gripper thread,
which applies the same grasp handling as in teleoperation.

The operator starts a run with a key after confirming the distance to the
first action. Inference runs only while the robot runs: each start clears the
old chunks and follows the first new one, because chunks predicted while the
arm stood still plan motion that did not happen. The robot holds its target when
the operator presses the hold key, when the newest chunk is older than 0.5~s,
when an action lies more than 100~mm from the current target, or when 180~s
have passed since the start; it stops when no chunk arrives for 3~s.
